\chapter*{STRESZCZENIE}
\addcontentsline{toc}{chapter}{\protect\hspace{0.45cm} STRESZCZENIE}

\noindent \textbf{Tytuł pracy w~języku polskim:}\\
Wpływ metod czyszczenia danych na jakość klasyfikatorów

\vspace{0.4cm}
\noindent \textbf{Tytuł pracy w~języku angielskim:}\\
Impact of Data Cleaning Methods on Classifier Performance

\vspace{0.4cm}
\noindent \textbf{Streszczenie:}\\
Niniejsza praca magisterska poświęcona jest empirycznej analizie wpływu metod wstępnego przetwarzania i~czyszczenia danych tabelarycznych na skuteczność oraz stabilność wybranych algorytmów klasyfikacji. Głównym problemem badawczym jest fakt, że poszczególne rodziny klasyfikatorów wykazują fundamentalnie odmienną wrażliwość na obecność zanieczyszczeń danych, co sprawia, że dobór i~kolejność transformacji preprocessing'owych mają bezpośrednie przełożenie na jakość uzyskiwanych modeli.

W~ramach pracy zdefiniowano i~zaimplementowano siedem komplementarnych strategii czyszczenia danych, obejmujących: punkt odniesienia bez interwencji (\textit{raw}), imputację medianą (\textit{fill}), imputację algorytmem $k$-najbliższych sąsiadów (\textit{fill\_knn}), filtrację wartości odstających regułą $3{,}0 \cdot \mathrm{IQR}$ z~następową imputacją medianą (\textit{remove\_fill}), standaryzację $z$-score po imputacji (\textit{fill\_norm}) oraz dwa pełne potoki wieloetapowe: klasyczny (\textit{all}) i~oparty na~$k$-NN (\textit{all\_knn}). Strategie zaimplementowano jako własne transformatory zgodne ze standardem \textsc{scikit-learn}, gwarantujące brak wycieku danych pomiędzy zbiorem uczącym a~testowym.

Badania empiryczne przeprowadzono na pięciu zbiorach danych z~różnych dziedzin: klinicznym (\textit{Zapalenia naczyń}, 848~próbek, 143~cechy), kardiologicznym (\textit{Serce}, 270$\times$13), diagnostycznym (\textit{Cukrzyca}, 768$\times$8), telekomunikacyjnym (\textit{Rezygnacje klientów}, 3333$\times$19) oraz finansowym (\textit{Ryzyko kredytowe}, 1000$\times$10). Każdy zbiór poddano kontrolowanej, trójskładnikowej degradacji jakości (losowe braki danych, ekstremalne wartości odstające z~mnożnikami $\times$20--80, szum kategoryczny) na pięciu poziomach od 10\% do~50\%, w~trzech niezależnych powtórzeniach losowych. Jako miarę ewaluacji przyjęto pole pod krzywą ROC ($\overline{\mathrm{AUC}}$) wyznaczone w~procedurze powtarzanej stratyfikowanej walidacji krzyżowej dla czterech rodzin klasyfikatorów: Lasów Losowych, XGBoost, Naiwnego Klasyfikatora Bayesa oraz Perceptronu Wielowarstwowego (MLP).

Kluczowe wyniki wykazały silną asymetrię wrażliwości badanych modeli. Sieć MLP bez czyszczenia przy $50\%$ uszkodzeń degraduje do poziomu klasyfikatora losowego ($\overline{\mathrm{AUC}} \approx 0{,}504$), natomiast zastosowanie pełnego potoku \textit{all} przywraca jej skuteczność do $0{,}642$ (zysk $\Delta\mathrm{AUC} = {+}0{,}138$). Naiwny Klasyfikator Bayesa zyskuje średnio $+0{,}10$ AUC po odfiltrowaniu wartości odstających. Modele zespołów drzew (Random Forest, XGBoost) wykazały wysoką naturalną odporność, a~czyszczenie przynosi im zysk głównie przy niskich poziomach uszkodzeń (10\%--30\%). Wykazano ponadto, że imputacja $k$-NN bez uprzedniej filtracji anomalii (\textit{fill\_knn}) jest metodą o~najniższej globalnej jakości spośród badanych, ponieważ nieskorygowane błędy grube zniekształcają metrykę odległości euklidesowej.

Istotność statystyczną wyników zweryfikowano dwuetapowo. Sparowane testy rangowe Wilcoxona potwierdziły istotną przewagę czyszczenia nad danymi surowymi w~77~spośród 120 przeprowadzonych porównań ($64{,}2\%$), w~tym 54-krotnie na poziomie $p < 0{,}001$. Globalny test Friedmana ($\chi^2_F = 544{,}06$, $p = 2{,}70 \times 10^{-114}$) i~procedura post-hoc Nemenyi z~Diagramem Różnic Krytycznych ($\mathrm{CD} = 0{,}300$) jednoznacznie wskazały potok \textit{all} jako statystycznie najlepszą strategię ($\bar{R} = 3{,}03$), istotnie lepszą od każdej pozostałej badanej metody.

\vspace{0.4cm}
\noindent \textbf{Słowa kluczowe:}\\
czyszczenie danych, uczenie maszynowe, klasyfikacja tabelaryczna, imputacja braków danych, wartości odstające, standaryzacja, walidacja krzyżowa, AUC, Random Forest, XGBoost, MLP, Gaussian Naive Bayes, test Friedmana, test Wilcoxona, diagram różnic krytycznych.

\vspace{0.8cm}
\noindent \textbf{Abstract:}\\
This master's thesis presents a~comprehensive empirical study of the impact of tabular data preprocessing and cleaning techniques on the predictive performance and stability of machine learning classifiers. The central research problem is that different classifier families exhibit fundamentally distinct sensitivity to data contamination, making the selection and ordering of preprocessing operations a~critical factor in model quality.

Seven complementary cleaning strategies were designed and implemented as \textsc{scikit-learn}-compatible transformers that prevent data leakage between training and test sets: a~no-intervention baseline (\textit{raw}), median imputation (\textit{fill}), $k$-Nearest Neighbor imputation (\textit{fill\_knn}), IQR-based outlier filtering followed by median imputation (\textit{remove\_fill}), imputation with $z$-score standardization (\textit{fill\_norm}), and two complete multi-step pipelines --- one based on median imputation (\textit{all}) and one based on $k$-NN imputation (\textit{all\_knn}).

Experiments were conducted on five benchmark datasets from diverse domains: clinical (\textit{Vasculitis}, 848~samples, 143~features), cardiology (\textit{Heart Disease}, 270$\times$13), diagnostics (\textit{Diabetes}, 768$\times$8), telecommunications (\textit{Customer Churn}, 3333$\times$19), and finance (\textit{German Credit}, 1000$\times$10). Each dataset was subjected to controlled three-component degradation (random missing values, extreme outliers with multipliers $\times$20--80, and categorical noise) at five contamination levels ranging from 10\% to 50\%, repeated three times with independent random seeds. The primary evaluation metric was the mean Area Under the ROC Curve ($\overline{\mathrm{AUC}}$) measured via repeated stratified cross-validation across four classifier families: Random Forest, XGBoost, Gaussian Naive Bayes, and Multilayer Perceptron (MLP).

Results revealed a~strong asymmetry in classifier sensitivity. Without cleaning, MLP degraded to a~random classifier at 50\% contamination ($\overline{\mathrm{AUC}} \approx 0.504$), while the \textit{all} pipeline restored its effectiveness to $0.642$ ($\Delta\mathrm{AUC} = {+}0.138$). Gaussian Naive Bayes gained an average of $+0.10$ AUC from outlier removal, while tree ensemble methods (Random Forest, XGBoost) exhibited high inherent robustness with preprocessing benefits concentrated at lower contamination levels (10\%--30\%). Notably, $k$-NN imputation without prior outlier filtering (\textit{fill\_knn}) ranked last globally, as uncorrected outliers distort the Euclidean distance metric and degrade imputation quality.

Statistical significance was verified in two stages. Paired Wilcoxon signed-rank tests confirmed a~significant advantage of cleaning over raw data in 77 out of 120 comparisons (64.2\%), including 54 at the $p < 0.001$ level. A~global Friedman test ($\chi^2_F = 544.06$, $p = 2.70 \times 10^{-114}$) followed by Nemenyi post-hoc analysis with Critical Difference Diagram ($\mathrm{CD} = 0.300$) unambiguously identified the \textit{all} pipeline as the statistically superior strategy ($\bar{R} = 3.03$), significantly outperforming every other evaluated method.

\vspace{0.4cm}
\noindent \textbf{Keywords:}\\
data cleaning, machine learning, tabular classification, missing value imputation, outlier detection, standardization, cross-validation, AUC, Random Forest, XGBoost, MLP, Gaussian Naive Bayes, Friedman test, Wilcoxon test, critical difference diagram.
